\chapter{Bilayer modular state and spectral resolvent}
\label{rm:chap:estado-modular-resolvente}

\section{The angular operator as common antecedent}

The constitutive operator of \cref{rm:ch:vacuum} admits a prior formulation
in which intensity and orientation occupy distinct coordinates of
the same object. Let \(R=R^*=R^{-1}\) be an involution with spectral multiplicities
one and one, and let
\begin{equation}
 P_+=\frac{I+R}{2},
 \qquad
 P_-=\frac{I-R}{2}
 \label{rm:eq:modular-projectors}
\end{equation}
their projectors. For
\begin{equation}
 x=A_{\rm rad}>0,
 \qquad
 y=C^*_{\rm rad},
 \qquad
 |y|<x,
 \label{rm:eq:modular-domain}
\end{equation}
we define the angular generator and its positive exponential by
\begin{equation}
 K_{\angle}=xI+yR,
 \qquad
 T_{x,y}=\ee^{-K_{\angle}}
 =q_+P_++q_-P_-,
 \label{rm:eq:modular-angular-operator}
\end{equation}
where
\begin{equation}
 q_+=\ee^{-x-y},
 \qquad
 q_-=\ee^{-x+y}.
 \label{rm:eq:modular-qpm}
\end{equation}
The condition \(|y|<x\) places both eigenvalues in \((0,1)\). The
determinant and the spectral quotient separate the two coordinates:
\begin{equation}
 D_A:=\det T_{x,y}=q_+q_-=\ee^{-2x},
 \qquad
 \varrho_{\rm or}:=\frac{q_-}{q_+}=\ee^{2y}.
 \label{rm:eq:modular-determinant-ratio}
\end{equation}
The first character is even under sheet exchange; the second records their orientation. Recovery is immediate:
\begin{equation}
 x=-\frac12\log D_A,
 \qquad
 y=\frac12\log\varrho_{\rm or}.
 \label{rm:eq:modular-recover-xy}
\end{equation}

The causal provenance of \(x\) and \(y\) remains linked to
state \(\mathrm{APP}\to\TRIT\to\TPK\). In particular,
\begin{equation}
 x=\frac{50\pi}{9}\,\mathfrak A_\beta,
 \qquad
 D_A=\exp\!\left(-\frac{100\pi}{9}\mathfrak A_\beta\right),
 \label{rm:eq:modular-alpha-determinant}
\end{equation}
where \(\mathfrak A_\beta=\alpha_{\HMT}\) denotes the terminal character of the
already constructed dodecaphase macro-object. The coordinate \(y\) subsequently comes
from the regional reader and preserves the orientation that the determinant eliminates.

\begin{proposition}[Canonical separation of intensity and orientation]
\label{rm:prop:modular-separation}
The map
\[
 (x,y)\longmapsto(D_A,\varrho_{\rm or})
 =(\ee^{-2x},\ee^{2y})
\]
is a diffeomorphism between the domain \(x>0\), \(|y|<x\), and the domain
\[
 0<D_A<1,
 \qquad
 D_A<\varrho_{\rm or}<D_A^{-1}.
\]
The sheet involution acts by
\(
(D_A,\varrho_{\rm or})\mapsto(D_A,\varrho_{\rm or}^{-1})
\).
\end{proposition}

\begin{proof}
The inverse formulas are \cref{rm:eq:modular-recover-xy}. The inequality
\(|y|<x\) is equivalent to
\(|\log\varrho_{\rm or}|<-\log D_A\), from which the stated interval is obtained.
Replacing \(y\) by \(-y\) preserves \(D_A\) and inverts the
ratio.
\end{proof}

\section{Normalization of the bilayer state}

The trace of the angular operator is
\begin{equation}
 \Tr T_{x,y}=2\ee^{-x}\cosh y.
 \label{rm:eq:modular-trace-T}
\end{equation}
Its normalization exactly eliminates the intensity \(x\) and preserves the
orientation:
\begin{equation}
 \boxed{
 \rho_y:=\frac{T_{x,y}}{\Tr T_{x,y}}
 =\frac{\ee^{-yR}}{2\cosh y}
 =p_+(y)P_++p_-(y)P_- ,}
 \label{rm:eq:modular-rho}
\end{equation}
with
\begin{equation}
 p_+(y)=\frac{\ee^{-y}}{2\cosh y},
 \qquad
 p_-(y)=\frac{\ee^{y}}{2\cosh y},
 \qquad
 p_++p_-=1.
 \label{rm:eq:modular-probabilities}
\end{equation}
The family \(\rho_y\) is faithful for every \(y\in\R\). Its orientation
observable satisfies
\begin{equation}
 \Tr(\rho_yR)=-\tanh y,
 \qquad
 \rho_{-y}=R_{\rm sh}\rho_yR_{\rm sh}^*,
 \label{rm:eq:modular-orientation-expectation}
\end{equation}
where \(R_{\rm sh}\) denotes any unitary operator permuting
\(P_+\) and \(P_-\). This implementation distinguishes the spectral involution
\(R\), which measures the leaf sign, from the unitary operator \(R_{\rm sh}\),
which exchanges the leaves.

\begin{definition}[Local modular holographic Hamiltonian]
\label{rm:def:holographic-modular-hamiltonian}
The modular Hamiltonian of the bilayer state is
\begin{equation}
 \boxed{
 K_y^{\rm mod}:=-\log\rho_y
 =\log(2\cosh y)I+yR.}
 \label{rm:eq:modular-Hamiltonian}
\end{equation}
The scalar component normalizes the state, and the \(yR\) component
preserves oriented memory.
\end{definition}

The adjective \emph{holographic} records the provenance of \(R\),
\(P_\pm\), and \(y\) from the enriched state and its two sheets. The equality
\cref{rm:eq:modular-Hamiltonian} is an identity of functional calculus; its
promotion to a generator of temporal dynamics requires specification of the
corresponding physical flow.

\section{Entropy, Parity, and Orientation Divergence}

\begin{theorem}[Exact informational geometry of the two sheets]
\label{rm:thm:modular-information-geometry}
The von Neumann entropy, parity, and relative entropy between the
two orientations satisfy
\begin{align}
 S(\rho_y)
 &:=-\Tr(\rho_y\log\rho_y)
 =\log(2\cosh y)-y\tanh y,
 \label{rm:eq:modular-entropy}\\
 S(\rho_y)&=S(\rho_{-y}),
 \label{rm:eq:modular-entropy-even}\\
 \boxed{
 D(\rho_y\Vert\rho_{-y})
 =2y\tanh y.}
 \label{rm:eq:modular-relative-entropy}
\end{align}
The divergence is positive for \(y\ne0\) and vanishes exactly on the
unoriented sheet \(y=0\).
\end{theorem}

\begin{proof}
From \cref{rm:eq:modular-Hamiltonian} and
\(\Tr(\rho_yR)=-\tanh y\) follows
\[
 S(\rho_y)=\Tr(\rho_yK_y^{\rm mod})
 =\log(2\cosh y)-y\tanh y.
\]
Both terms are even. Moreover,
\[
 \log\rho_y-\log\rho_{-y}=-2yR,
\]
and, consequently,
\[
 D(\rho_y\Vert\rho_{-y})
 =-2y\Tr(\rho_yR)=2y\tanh y.
\]
The product \(y\tanh y\) is positive for \(y\ne0\).
\end{proof}

Entropy quantifies the even distribution of the sheets, whereas
divergence quantifies the informational separation of their orientations.
The derivatives
\begin{equation}
 \frac{\dd}{\dd y}S(\rho_y)=-y\operatorname{sech}^2y,
 \qquad
 \frac{\dd^2}{\dd y^2}\log(2\cosh y)
 =\operatorname{sech}^2y
 \label{rm:eq:modular-entropy-derivatives}
\end{equation}
identify \(S(\rho_0)=\log2\) as a maximum and provide the Fisher information of the exponential family. Near the symmetric sheet,
\begin{equation}
 D(\rho_y\Vert\rho_{-y})
 =2y^2-\frac23y^4+O(y^6).
 \label{rm:eq:modular-relative-expansion}
\end{equation}
The first contribution is quadratic: a small oriented memory
has a second-order informational cost.

The trace distance between orientations is also obtained in
closed form:
\begin{equation}
 \frac12\lVert\rho_y-\rho_{-y}\rVert_1=|\tanh y|.
 \label{rm:eq:modular-trace-distance}
\end{equation}
Since \(|y|\geq|\tanh y|\), the identity
\cref{rm:eq:modular-relative-entropy} implies the control
\begin{equation}
 D(\rho_y\Vert\rho_{-y})
 \geq 2\left(\frac12
 \lVert\rho_y-\rho_{-y}\rVert_1\right)^2,
 \label{rm:eq:modular-pinsker-exact-control}
\end{equation}
consistent with Pinsker's inequality. Thus, the oriented quantity simultaneously has
spectral, entropic, and metric readings.

\section{Normalized resolvent of the angular macro-object}

Let
\begin{equation}
 \kappa_\beta:=4\pi\mathfrak A_\beta>0.
 \label{rm:eq:modular-kappa-beta}
\end{equation}
We define the meromorphic matrix-valued function
\begin{equation}
 \boxed{
 \mathfrak M_\beta(z)
 :=\kappa_\beta^{-1}(zI-T_{x,y})^{-1},
 \qquad
 z\in\C\setminus\{q_+,q_-\}.}
 \label{rm:eq:modular-resolvent}
\end{equation}
Its spectral decomposition is
\begin{equation}
 \mathfrak M_\beta(z)
 =\frac1{\kappa_\beta}
 \left(\frac{P_+}{z-q_+}+\frac{P_-}{z-q_-}\right).
 \label{rm:eq:modular-resolvent-spectral}
\end{equation}
The resolvent unites in a single object the coupling scale, the angular
poles, the sheet projectors, the determinant, and all subsequent functional
calculus.

\begin{theorem}[Injectivity by the first two moments]
\label{rm:thm:modular-resolvent-injective}
The assignment
\[
 (\mathfrak A_\beta,T_{x,y})
 \longmapsto\mathfrak M_\beta
\]
is injective for \(\mathfrak A_\beta>0\). specifically, if
\begin{align}
 B_0&:=\lim_{|z|\to\infty}z\mathfrak M_\beta(z),
 \label{rm:eq:modular-B0}\\
 B_1&:=\lim_{|z|\to\infty}z^2
 \left(\mathfrak M_\beta(z)-z^{-1}B_0\right),
 \label{rm:eq:modular-B1}
\end{align}
then
\begin{equation}
 B_0=\kappa_\beta^{-1}I,
 \qquad
 \kappa_\beta=\tau_2(B_0)^{-1},
 \qquad
 T_{x,y}=\kappa_\beta B_1,
 \label{rm:eq:modular-reconstruct-kappa-T}
\end{equation}
where \(\tau_2=\Tr/2\). Therefore, the asymptotics recovers
\(\mathfrak A_\beta=\kappa_\beta/(4\pi)\), and the next moment recovers
the complete angular operator.
\end{theorem}

\begin{proof}
For \(|z|>\lVert T_{x,y}\rVert\), the Neumann series gives
\begin{equation}
 \mathfrak M_\beta(z)
 =\frac1{\kappa_\beta}
 \sum_{n\geq0}\frac{T_{x,y}^{n}}{z^{n+1}}.
 \label{rm:eq:modular-Laurent}
\end{equation}
The coefficients of \(z^{-1}\) and \(z^{-2}\) are precisely \(B_0\) and
\(B_1\). The equalities \cref{rm:eq:modular-reconstruct-kappa-T} follow upon
taking the normalized trace and multiplying by \(\kappa_\beta\). Two pairs with
the same resolvent have the same coefficients and therefore coincide.
\end{proof}

If \(q_+\ne q_-\), the spectral data are also recovered from the poles
and residues:
\begin{equation}
 \operatorname*{Res}_{z=q_\pm}\mathfrak M_\beta(z)
 =\kappa_\beta^{-1}P_\pm,
 \qquad
 \det\mathfrak M_\beta(z)
 =\frac{\kappa_\beta^{-2}}
 {(z-q_+)(z-q_-)}.
 \label{rm:eq:modular-resolvent-residues}
\end{equation}
in particular, the scalar denominator contains
\(q_++q_-\) and \(q_+q_-=D_A\), and
\begin{equation}
 P_+=\frac{T-q_-I}{q_+-q_-},
 \qquad
 P_-=\frac{T-q_+I}{q_--q_+}.
 \label{rm:eq:modular-projector-recovery}
\end{equation}
The case \(y=0\) produces a unique simple pole of spectral multiplicity
two: \(T=\ee^{-x}I\). The resolvent then preserves the intensity and the
scalar operator, whereas the choice of projectors has no intrinsic
content because the orientation is exactly zero.

\begin{corollary}[Calculation functional from the resolvente]
\label{rm:cor:modular-functional-calculus}
For every function \(f\) holomorphic in a neighborhood of
\(\{q_+,q_-\}\),
\begin{equation}
 f(T)=\frac{\kappa_\beta}{2\pi\ii}
 \int_\Gamma f(z)\mathfrak M_\beta(z)\,\dd z,
 \label{rm:eq:modular-Dunford}
\end{equation}
where \(\Gamma\) encircles both poles. In particular, the vacuum response
\[
 \mathcal V=(I-T^{90})(I-T^{120})^{-1}
\]
is reconstructed from \(\mathfrak M_\beta\) without incorporating additional
scalar data.
\end{corollary}

\begin{proof}
The identity is the formula of Dunford--Riesz aplicada to
\((zI-T)^{-1}=\kappa_\beta\mathfrak M_\beta(z)\). The function
\(f(z)=(1-z^{90})/(1-z^{120})\) is holomorphic over a neighborhood of the spectrum,
since \(0<q_\pm<1\).
\end{proof}

\section{Compatibility with nonadic refinement}

The preceding information is prolonged through the tower
\begin{equation}
 \mathcal A_m=M_2(\C)\otimes M_{9^m}(\C),
 \qquad
 \iota_m(a)=a\otimes I_9.
 \label{rm:eq:modular-UHF-tower}
\end{equation}
Let
\begin{equation}
 T_m=T\otimes I_{9^m},
 \quad
 R_m=R\otimes I_{9^m},
 \quad
 \mathfrak M_{\beta,m}(z)
 =\kappa_\beta^{-1}(zI-T_m)^{-1},
 \label{rm:eq:modular-UHF-objects}
\end{equation}
and let
\(E_m=\operatorname{id}_{M_2\otimes M_{9^m}}\otimes\tau_9\) be the
tracial expectation from \(\mathcal A_{m+1}\) to \(\mathcal A_m\).

\begin{theorem}[Resolvent and modular naturality]
\label{rm:thm:modular-UHF-naturality}
for every \(m\geq0\) and every \(z\) of the resolvent set,
\begin{align}
 E_m(T_{m+1})&=T_m,
 \label{rm:eq:modular-UHF-T}\\
 E_m\!\left(\mathfrak M_{\beta,m+1}(z)\right)
 &=\mathfrak M_{\beta,m}(z),
 \label{rm:eq:modular-UHF-resolvent}\\
 E_m\!\left(f(T_{m+1})\right)&=f(T_m)
 \label{rm:eq:modular-UHF-functional}
\end{align}
for every continuous function \(f\) on the spectrum. Likewise, the state functionals
\begin{equation}
 \omega_{y,m}(a)
 :=(\Tr_2\otimes\tau_{9^m})
 \bigl((\rho_y\otimes I_{9^m})a\bigr)
 \label{rm:eq:modular-UHF-state}
\end{equation}
are compatible with the inclusions and define a state \(\omega_{y,\infty}\)
on the UHF limit.
\end{theorem}

\begin{proof}
All objects of \cref{rm:eq:modular-UHF-objects} have the form
\(a\otimes I_{9^m}\). The tracial expectation eliminates the last identity
factor and preserves \(a\). The same observation proves
\cref{rm:eq:modular-UHF-functional}; compatibility of the states is
obtained by evaluating \(\omega_{y,m+1}(a\otimes I_9)\) and using
\(\tau_9(I_9)=1\).
\end{proof}

The local entropy relative to the trace
\(\Tr_2\otimes\tau_{9^m}\) remains equal to
\(S(\rho_y)\), and the divergence between orientations remains equal to
\(2y\tanh y\). Refinement adds resolution that is nonadic and preserves the
bilayer content. In this manner, the construction lives in the UHF tower and
in its tracial closure, rather than depending on an accidental presentation
by matrices \(2\times2\).

\section{Exactness of the modular resolvent, normalization obstructions, and separation of the thermodynamic reading}

\begin{proofstatusbox}
The decomposition of \(T\), the state \(\rho_y\), the modular Hamiltonian,
entropy, relative divergence, injectivity of the resolvent, and
UHF naturality are exact results once
\((\mathfrak A_\beta,C^*,R)\) has been fixed. The assembly of these identities in the
resolvent \(\mathfrak M_\beta\) constitutes a new formalization of the
preexisting authorial architecture. The reading of
\(D(\rho_y\Vert\rho_{-y})\) as physical work requires temperature,
protocol, and dynamics; until these are fixed, its status is that of an exact
informational measure. Realizations through Barbero, CKM, or Higgs are
subsequent readers and preserve their respective intertwiners.
\end{proofstatusbox}

\begin{falsifierbox}
The block is refuted if the normalization of \(T\) retains a dependence on \(x\), if \(K_y^{\rm mod}\) does not exponentiate to \(\rho_y\), if the relative entropy differs from \(2y\tanh y\), if two distinct pairs \((\mathfrak A_\beta,T)\) have the same resolvent, or if any expectation \(E_m\) alters the functional calculation. Recovery of the projectors in \(y=0\) additionally constitutes a negative control: any sheet selector that survived that degeneration would have been introduced externally.
\end{falsifierbox}
